%% kapa: exam=1 round=35 year=2024 period=1 form=A subject=civil_law subject_ko=민법 seq=18 number=18
%% source: Q-Net 감정평가사 기출문제 (cst00309) artlSeq=5213229 file=raw/1008/5213229/1교시 문제지 원본.pdf
%% conversion: kordoc 4.24.51 -> tools/gen1.py (auto); review=proofread-sonnet; flags=none
\begin{question}{18}
소멸시효의 기산점이 잘못 연결된 것은? (다툼이 있으면 판례에 따름)
\begin{choices}
  \item 불확정기한부 채권 - 기한이 객관적으로 도래한 때
  \item 부당이득반환청구권 - 기한의 도래를 안 때
  \item 정지조건부 권리 - 조건이 성취된 때
  \item 부작위를 목적으로 하는 채권 - 위반행위를 한 때
  \item 선택채권 - 선택권을 행사할 수 있을 때
\end{choices}
\end{question}
\answer{2}
